% ============================================================
% FRANÇAIS — IMPARFAIT DE L'INDICATIF
% Version graphique LaTeX/TikZ
% Compilation : XeLaTeX
% ============================================================
\documentclass[a4paper,10pt]{article}

\usepackage[a4paper,top=9mm,bottom=8mm,left=11mm,right=11mm]{geometry}
\usepackage{fontspec}
\usepackage[french]{babel}
\usepackage{microtype}
\usepackage{tikz}
\usetikzlibrary{arrows.meta,calc,positioning}
\usepackage[most]{tcolorbox}
\usepackage{tabularx}
\usepackage{array}
\usepackage{xcolor}
\usepackage[normalem]{ulem}

\setmainfont{Noto Sans}
\setsansfont{Noto Sans}

% ------------------------------------------------------------
% PALETTE — PASSÉ = BLEU
% ------------------------------------------------------------
\definecolor{Blue}{HTML}{356FA8}
\definecolor{BlueDark}{HTML}{244D78}
\definecolor{BlueLight}{HTML}{EAF2F9}
\definecolor{BlueLine}{HTML}{9BBBD7}
\definecolor{Red}{HTML}{D52F35}
\definecolor{RedLight}{HTML}{FCEBED}
\definecolor{Gold}{HTML}{B27A00}
\definecolor{GoldLight}{HTML}{FFF3D7}
\definecolor{Ink}{HTML}{172A32}
\definecolor{Grey}{HTML}{64757D}
\definecolor{Border}{HTML}{B9D3E3}
\definecolor{White}{HTML}{FFFFFF}

\pagestyle{empty}
\setlength{\parindent}{0pt}
\setlength{\parskip}{0pt}
\setlength{\tabcolsep}{3pt}

% ------------------------------------------------------------
% CADRE EXTÉRIEUR
% ------------------------------------------------------------
\newcommand{\PageFrame}{%
\begin{tikzpicture}[remember picture,overlay]
  \coordinate (A) at ([xshift=8mm,yshift=-8mm]current page.north west);
  \coordinate (B) at ([xshift=-8mm,yshift=-8mm]current page.north east);
  \coordinate (C) at ([xshift=-8mm,yshift=8mm]current page.south east);
  \coordinate (D) at ([xshift=8mm,yshift=8mm]current page.south west);
  \def\r{1.7mm}

  \draw[Border,line width=.65pt]
    (A)
    -- ([xshift=-\r]B)
    arc[start angle=90,end angle=0,radius=\r]
    -- ([yshift=\r]C)
    arc[start angle=0,end angle=-90,radius=\r]
    -- (D) -- cycle;

  \draw[Blue,line width=3.2mm] (A) -- (D);
\end{tikzpicture}%
}

% ------------------------------------------------------------
% BOÎTES
% ------------------------------------------------------------
\tcbset{
  enhanced,
  boxrule=.6pt,
  arc=2.8mm,
  left=4mm,right=4mm,top=2mm,bottom=2mm,
  coltext=Ink
}

\newtcolorbox{timelinebox}{colback=White,colframe=BlueLine}
\newtcolorbox{greenbox}{colback=BlueLight,colframe=BlueLine}
\newtcolorbox{bluebox}{colback=BlueLight,colframe=Border}
\newtcolorbox{redbox}{colback=RedLight,colframe=Border}
\newtcolorbox{whitebox}{colback=White,colframe=Border}
\newtcolorbox{goldbox}{colback=GoldLight,colframe=GoldLight,boxrule=0pt,arc=2.5mm,left=2.7mm,right=2.7mm,top=1.25mm,bottom=1.25mm}

% ------------------------------------------------------------
% PETITES ÉTIQUETTES
% ------------------------------------------------------------
\newcommand{\pill}[2][BlueLight]{%
  \tikz[baseline=(X.base)]{\node[fill=#1,rounded corners=2pt,inner xsep=4.2pt,inner ysep=1.6pt,font=\small\bfseries,text=Ink] (X) {#2};}%
}
\newcommand{\sectiontitle}[2]{{\large\bfseries\textcolor{#1}{#2}}}
\newcommand{\signalpill}[1]{%
  \tikz[baseline=(X.base)]{\node[fill=GoldLight,rounded corners=2pt,inner xsep=3.7pt,inner ysep=1.1pt,font=\scriptsize\bfseries,text=Ink] (X) {#1};}%
}

% ------------------------------------------------------------
% ------------------------------------------------------------
% LIGNE DU TEMPS — IMPARFAIT FRANÇAIS
% ------------------------------------------------------------
\newcommand{\Timeline}{%
\begin{tikzpicture}[x=1cm,y=1cm]
  \draw[-{Latex[length=3mm]},very thick,Ink] (0,0) -- (15.3,0);

  % repère du passé : l'action se déroule sur une durée passée
  \draw[-{Latex[length=3mm,width=2mm]},very thick,Blue]
    (6.6,1.10) -- (6.6,.10);
  \node[font=\bfseries\small,text=Blue] at (6.6,1.32) {MOMENT PASSÉ};

  % zone de durée / arrière-plan
  \draw[Blue,line width=3.2pt] (3.0,0) -- (7.8,0);
  \fill[Blue] (3.0,0) circle (2.1pt);
  \fill[Blue] (7.8,0) circle (2.1pt);

  \node[font=\scriptsize\bfseries,text=Grey] at (2.2,-.48) {PASSÉ};
  \node[font=\scriptsize\bfseries,text=Blue] at (6.6,-.48) {IMPARFAIT};
  \node[font=\scriptsize\bfseries,text=Grey] at (13.3,-.48) {FUTUR};

  \node[fill=BlueLight,rounded corners=2pt,inner xsep=5pt,inner ysep=2pt,font=\scriptsize\bfseries,text=BlueDark]
    at (5.0,.72) {durée / habitude dans le passé};
\end{tikzpicture}%
}

\begin{document}
\PageFrame

% ------------------------------------------------------------
% TITRE
% ------------------------------------------------------------
\vspace*{1mm}
\begin{center}
  {\fontsize{25}{27}\selectfont\bfseries\textcolor{Ink}{IMPARFAIT}}

  \vspace{0.75mm}
  {\large\textcolor{Grey}{L'imparfait de l'indicatif}}

  \vspace{3mm}

  \begin{tcolorbox}[
    width=.88\textwidth,
    colback=BlueLight,
    colframe=BlueLight,
    boxrule=0pt,
    arc=3mm,
    halign=center,
    top=2.2mm,bottom=2.2mm
  ]
    {\large\bfseries\textcolor{BlueDark}{NOUS AU PRÉSENT + TERMINAISONS DE L'IMPARFAIT}}
    \quad
    {\small\textit{ex. : parlons $\rightarrow$ parl- + ais, ais, ait, ions, iez, aient}}
  \end{tcolorbox}
\end{center}

\vspace{1mm}

% ------------------------------------------------------------
% TIMELINE
% ------------------------------------------------------------
\begin{timelinebox}
  \sectiontitle{Blue}{La ligne du temps}

  \vspace{0.75mm}
  \begin{center}
    \Timeline
  \end{center}

  \vspace{-1mm}
  \begin{center}
    \small
    L'\textbf{imparfait} situe une action ou un état dans le \textbf{passé},
    en insistant souvent sur sa \textbf{durée, son habitude ou son contexte}.
  \end{center}
\end{timelinebox}

\vspace{1.35mm}

% ------------------------------------------------------------
% ------------------------------------------------------------
% FORMATION
% ------------------------------------------------------------
\begin{greenbox}
  \sectiontitle{Blue}{Formation — exemple : parler}

  \vspace{0.75mm}
  \begin{center}
    \begin{tabular}{@{}c@{\qquad}c@{\qquad}c@{\qquad}c@{\qquad}c@{\qquad}c@{}}
      je & tu & il / elle / on & nous & vous & ils / elles\\[-.2mm]
      \textcolor{Blue}{parlais} & \textcolor{Blue}{parlais} & \textcolor{Blue}{parlait} &
      \textcolor{Blue}{parlions} & \textcolor{Blue}{parliez} & \textcolor{Blue}{parlaient}
    \end{tabular}
  \end{center}

  \vspace{0.5mm}
  \begin{center}
    \small\textcolor{Grey}{On part en général de la forme \textbf{nous} au présent, puis on ajoute :}
    \quad
    \textbf{-ais \; -ais \; -ait \; -ions \; -iez \; -aient}
  \end{center}
\end{greenbox}

\vspace{1.35mm}

% ------------------------------------------------------------
% ------------------------------------------------------------
% USES
% ------------------------------------------------------------
\begin{greenbox}
  \sectiontitle{Blue}{Quand utilise-t-on l'imparfait ?}

  \vspace{0.75mm}
  \begin{tabularx}{\textwidth}{@{}p{.31\textwidth} X@{}}
    \pill{HABITUDE PASSÉE}
      & Action répétée ou habituelle dans le passé.\newline
        \textit{Quand j'étais enfant, je jouais dehors.}\\[0.7mm]
    \pill{DURÉE / ACTION EN COURS}
      & Action qui était en train de se dérouler à un moment passé.\newline
        \textit{Je lisais quand tu es arrivé.}\\[0.7mm]
    \pill{DESCRIPTION / CONTEXTE}
      & Cadre, décor, état ou situation dans un récit au passé.\newline
        \textit{Il faisait froid et la rue était déserte.}\\[0.7mm]
    \pill{ÉTAT / SITUATION DURABLE}
      & Situation qui durait ou caractérisait une période passée.\newline
        \textit{Elle habitait alors à Albi.}
  \end{tabularx}
\end{greenbox}

\vspace{1.35mm}

% ------------------------------------------------------------
% ------------------------------------------------------------
% EXAMPLES
% ------------------------------------------------------------
\begin{bluebox}
  \sectiontitle{Blue}{Exemples}

  \vspace{0.75mm}
  \begin{tabularx}{\textwidth}{@{}X@{\hspace{7mm}}X@{}}
    \textbf{Je travaillais tous les samedis.}
    & \textbf{Il pleuvait ce matin-là.}\\[-.5mm]
    \textit{Habitude passée.}
    & \textit{Contexte / description.}\\[0.7mm]
    \textbf{Nous marchions dans la forêt.}
    & \textbf{Elle habitait dans le Tarn.}\\[-.5mm]
    \textit{Action en cours dans le passé.}
    & \textit{Situation durable dans le passé.}
  \end{tabularx}
\end{bluebox}

\vspace{1.35mm}

% ------------------------------------------------------------
% ------------------------------------------------------------
% POINT DE VIGILANCE
% ------------------------------------------------------------
\begin{redbox}
  \textcolor{Red}{\bfseries !\quad À retenir :}
  l'imparfait présente souvent le passé comme \textbf{en cours, habituel ou descriptif}.
  Il s'oppose notamment au passé composé, qui présente plus facilement un fait comme \textbf{achevé ou ponctuel}.
\end{redbox}

\vspace{1.35mm}

% ------------------------------------------------------------
% ------------------------------------------------------------
% MOTS REPÈRES
% ------------------------------------------------------------
\begin{goldbox}
  \begin{center}
    {\normalsize\bfseries\textcolor{Gold}{MOTS REPÈRES}}
    \hspace{3mm}
    \signalpill{souvent}\hspace{1mm}
    \signalpill{toujours}\hspace{1mm}
    \signalpill{habituellement}\hspace{1mm}
    \signalpill{autrefois}\hspace{1mm}
    \signalpill{à l'époque}

    \vspace{0.75mm}

    \signalpill{chaque jour}\hspace{1mm}
    \signalpill{tous les jours}\hspace{1mm}
    \signalpill{quand j'étais...}\hspace{1mm}
    \signalpill{pendant que}
  \end{center}
\end{goldbox}

\end{document}
